\subsection{Directed families; support of a measure}

PROPOSITION 5. --- a) Let H be an increasing directed set of functions \(\geqslant 0\) that are lower semi-continuous on every compact subset of T. Then

\[
\mu^ {\bullet} \left(\sup _ {h \in \mathrm{H}} h\right) = \sup _ {h \in \mathrm{H}} \mu^ {\bullet} (h).\tag{5}
\]

\begin{enumerate}
\def\labelenumi{\alph{enumi})}
\setcounter{enumi}{1}
\tightlist
\item
  Let H be a decreasing directed set of functions \(\geqslant0\) that are upper semi-continuous on every compact subset of T. If there exists in H a function \(h_{0}\) such that \(\mu^{\bullet}(h_{0})<+\infty\) , then
\end{enumerate}

\[
\mu^ {\bullet} \left(\inf _ {h \in \mathrm{H}} h\right) = \inf _ {h \in \mathrm{H}} \mu^ {\bullet} (h).\tag{6}
\]

For every compact set \(\mathbf{K} \subset \mathbf{T}\) , we have in case \(a\) )

\[
\mu^ {\bullet} \Big (\sup _ {h \in \mathrm{H}} h \varphi_ {\mathrm{K}} \Big) = \mu_ {\mathrm{K}} ^ {\bullet} \Big (\sup _ {h \in \mathrm{H}} h _ {\mathrm{K}} \Big) = \sup _ {h \in \mathrm{H}} \mu_ {\mathrm{K}} ^ {\bullet} (h _ {\mathrm{K}}) = \sup _ {h \in \mathrm{H}} \mu^ {\bullet} (h \varphi_ {\mathrm{K}}),
\]

and in case \(b\)

\[
\mu^ {\bullet} \Big (\inf _ {h \in \mathrm{H}} h \varphi_ {\mathrm{K}} \Big) = \mu_ {\mathrm{K}} ^ {\bullet} \Big (\inf _ {h \in \mathrm{H}} h _ {\mathrm{K}} \Big) = \inf _ {h \in \mathrm{H}} \mu_ {\mathrm{K}} ^ {\bullet} (h _ {\mathrm{K}}) = \inf _ {h \in \mathrm{H}} \mu^ {\bullet} (h \varphi_ {\mathrm{K}}),
\]

by Prop. 2 of No.~2, and Prop. 8 of Ch. V, §1, No.~2. Case a) follows at once, by passage to the supremum with respect to K (No.~2, Prop. 2). To treat case b), denote by \(\varepsilon\) a number \(>0\) , and choose a compact set K such that \(\mu^{\bullet}(h_0\varphi_{\mathbf{K}}) \geqslant \mu^{\bullet}(h_0) - \varepsilon\) . We then have (No.~5, Prop. 4) \(\mu^{\bullet}(h_0\varphi_{\mathbf{C}_{\mathbf{K}}}) \leqslant \varepsilon\) ; for every function \(h \in \mathbf{H}\) that is \(\leqslant h_0\) , we therefore have \(\mu^{\bullet}(h\varphi_{\mathbf{C}_{\mathbf{K}}}) \leqslant \varepsilon\) , and finally \(\mu^{\bullet}(h\varphi_{\mathbf{K}}) \geqslant \mu^{\bullet}(h) - \varepsilon\) by Prop. 4 of No.~5. Therefore

\[
\mu^ {\bullet} \left(\inf _ {h \in \mathrm{H}} h\right) \geqslant \mu^ {\bullet} \left(\inf _ {h \in \mathrm{H}} h \varphi_ {\mathrm{K}}\right) = \inf _ {h \in \mathrm{H}, h \leqslant h _ {0}} \mu^ {\bullet} (h \varphi_ {\mathrm{K}}) \geqslant \inf _ {h \in \mathrm{H}, h \leqslant h _ {0}} \mu^ {\bullet} (h) - \varepsilon .
\]

Consequently the left side of (6) is \(\geqslant\) the right side; the reverse inequality being obvious, the proposition is established.

COROLLARY. --- a) Let \((\mathrm{U}_{\alpha})_{\alpha \in \mathrm{I}}\) be an increasing directed family of open subsets of T, with union U. Then \(\mu^{\bullet}(\mathrm{U}) = \sup_{\alpha \in \mathrm{I}} \mu^{\bullet}(\mathrm{U}_{\alpha})\) .

\begin{enumerate}
\def\labelenumi{\alph{enumi})}
\setcounter{enumi}{1}
\tightlist
\item
  Let \((\mathbf{F}_{\alpha})_{\alpha \in \mathbf{I}}\) be a decreasing directed family of closed subsets of \(\mathbf{T}\) , with intersection \(\mathbf{F}\) . If there exists an \(\alpha \in \mathbf{I}\) such that \(\mu^{\bullet}(\mathbf{F}_{\alpha})\) is finite, then \(\mu^{\bullet}(\mathbf{F}) = \inf_{\alpha \in \mathbf{I}} \mu^{\bullet}(\mathbf{F}_{\alpha})\) .
\end{enumerate}

By the preceding corollary, there exists a largest locally negligible open set; this justifies the following definition:

DEFINITION 8. --- The support of a measure \(\mu\) on T is defined to be the complement of the largest locally \(\mu\) -negligible open set in T.

The support of \(\mu\) is denoted \(\operatorname{Supp}(\mu)\) .

Remarks. --- 1) If \(\mu\) is a complex measure, the support of \(\mu\) is defined to be the support of the positive measure \(|\mu|\) ; it is again the complement of the largest locally \(\mu\) -negligible open set.

\begin{enumerate}
\def\labelenumi{\arabic{enumi})}
\setcounter{enumi}{1}
\tightlist
\item
  Let us show that the measures introduced in Example 2 of No.~3 are measures with compact support in T. Let \(\mu\) be a positive measure on T whose support is a compact set K, and let \(\nu\) be the measure defined by \(\mu_{\mathbf{K}}\) (in the sense of No.~3). Let \(f \in \mathcal{F}_{+}(\mathbf{T})\) ; then
\end{enumerate}

\[
\nu^ {\bullet} (f) = \mu_ {K} ^ {\bullet} (f _ {K}) \quad (\text { No.   3,   formula   (3)) }.
\]

The encumbrance \(\mu^{\bullet}\) being concentrated on K, we also have

\[
\mu^ {\bullet} (f) = \mu^ {\bullet} (f \varphi_ {K}) = \mu^ {\bullet} ((f _ {K}) ^ {0}) = \mu_ {K} ^ {\bullet} (f _ {K}),
\]

whence \(\mu^{\bullet} = \nu^{\bullet}\) , and finally \(\mu = \nu\) . Conversely, if K is a compact set in T and \(\lambda\) a measure on K, and if \(\mu\) is the measure on T defined by \(\lambda\) , then \(\mu^{\bullet}(\mathbb{C}K) = 0\) (No.~3, formula (3)); consequently, the support of \(\mu\) is contained in K, hence is compact.

